\usepackage{tikz}
\usetikzlibrary{arrows.meta,positioning,calc,fit,backgrounds}
\usepackage{pgfplots}
\pgfplotsset{compat=1.18}
\usepackage{tabularx}
\graphicspath{{figures/}{../figures/}}
\newcommand{\protset}{\mathcal{P}}
\newcommand{\risk}{\mathcal{R}}
\newcommand{\E}{\mathbb{E}}
\newcommand{\Prob}{\mathbb{P}}
\newcommand{\bits}{\mathcal{B}}
\newcommand{\code}[1]{\texttt{#1}}
\newcommand{\clip}{\operatorname{clip}}
\definecolor{nsnavy}{HTML}{18354A}
\definecolor{nsteal}{HTML}{087E8B}
\definecolor{nsgold}{HTML}{BC771B}
\definecolor{nspale}{HTML}{EFF4F6}
% Draft markers (15 November 2026): red text is provisional evidence or a result still being computed.
\newcommand{\pending}[1]{{\color{red}#1}}
\newcommand{\waiting}[1]{{\color{red}\textbf{[Pending:} #1\textbf{]}}}
% Cross-references: the whole "Figure 3" / "Section 5" / "Eq. (4)" is the link, and links are boxed (hyperref's default
% look: red for internal links, green for citations, cyan for URLs) instead of coloured.
\newcommand{\wref}[2]{\hyperref[#2]{#1~\ref*{#2}}}
\newcommand{\weqref}[2]{\hyperref[#2]{#1~\textup{(\ref*{#2})}}}
\AtBeginDocument{\hypersetup{colorlinks=false,pdfborder={0 0 0.8},linkbordercolor={1 0 0},citebordercolor={0 0.6 0},urlbordercolor={0 0.6 0.9}}}
